\documentclass[12pt]{article}
\usepackage[utf8]{inputenc}
\usepackage[T1]{fontenc}
\usepackage[spanish]{babel}
\usepackage[a4paper,margin=2.5cm]{geometry}
\usepackage{amsmath}
\usepackage{tikz}
\usepackage{listings}
\usepackage{xcolor}
\usetikzlibrary{automata,positioning,arrows.meta}

\definecolor{codegray}{rgb}{0.95,0.95,0.95}
\lstset{
    basicstyle=\ttfamily\small,
    backgroundcolor=\color{codegray},
    frame=single,
    breaklines=true,
    columns=fullflexible
}

\title{Explicacion de la Implementacion de \texttt{uart\_tx}}
\author{}
\date{}

\begin{document}

\maketitle

\section{Objetivo del modulo}

El archivo \texttt{src/uart\_tx.v} implementa un transmisor UART sencillo. Su trabajo es tomar un byte de entrada (\texttt{tx\_data}) y enviarlo por la linea serie \texttt{tx} usando el formato clasico:

\begin{itemize}
    \item 1 bit de inicio (\texttt{start bit}) con valor \texttt{0}
    \item 8 bits de datos
    \item 1 bit de parada (\texttt{stop bit}) con valor \texttt{1}
\end{itemize}

Mientras se transmite el byte, la salida \texttt{tx\_busy} indica que el transmisor esta ocupado.

\section{Interfaz del modulo}

\begin{lstlisting}[language=Verilog]
module uart_tx #(
    parameter CLK_FREQ  = 27_000_000,
    parameter BAUD_RATE = 115200
)(
    input  wire       clk,
    input  wire       rst_n,
    input  wire       tx_en,
    input  wire [7:0] tx_data,
    output reg        tx,
    output reg        tx_busy
);
\end{lstlisting}

\subsection*{Entradas}

\begin{itemize}
    \item \texttt{clk}: reloj principal del sistema.
    \item \texttt{rst\_n}: reset activo en bajo. Cuando vale \texttt{0}, el modulo vuelve a su estado inicial.
    \item \texttt{tx\_en}: pulso o habilitacion que ordena iniciar una transmision.
    \item \texttt{tx\_data}: byte que se desea transmitir.
\end{itemize}

\subsection*{Salidas}

\begin{itemize}
    \item \texttt{tx}: linea serie de salida.
    \item \texttt{tx\_busy}: indica si el modulo esta enviando un dato.
\end{itemize}

\section{Parametros}

\begin{lstlisting}[language=Verilog]
localparam CLOCKS_PER_BIT = CLK_FREQ / BAUD_RATE;
\end{lstlisting}

Esta constante define cuantos ciclos de reloj debe durar cada bit en la UART. Por ejemplo, si:

\[
\texttt{CLK\_FREQ} = 27\,000\,000
\qquad\text{y}\qquad
\texttt{BAUD\_RATE} = 115200
\]

entonces:

\[
\texttt{CLOCKS\_PER\_BIT} \approx \frac{27\,000\,000}{115200} = 234
\]

Eso significa que cada bit se mantiene estable durante 234 ciclos de \texttt{clk}.

\section{Estados de la maquina}

\begin{lstlisting}[language=Verilog]
localparam s_IDLE  = 2'b00;
localparam s_START = 2'b01;
localparam s_DATA  = 2'b10;
localparam s_STOP  = 2'b11;
\end{lstlisting}

El modulo esta organizado como una maquina de estados finitos con cuatro etapas:

\begin{itemize}
    \item \texttt{s\_IDLE}: espera un nuevo dato para transmitir.
    \item \texttt{s\_START}: envia el bit de inicio.
    \item \texttt{s\_DATA}: envia los 8 bits del byte.
    \item \texttt{s\_STOP}: envia el bit de parada.
\end{itemize}

\begin{figure}[h!]
\centering
{\shorthandoff{<>}
\begin{tikzpicture}[
    >=Latex,
    node distance=30mm,
    on grid,
    auto,
    every state/.style={draw, thick, minimum size=11mm, font=\small}
]
    \node[state, initial] (idle) {\texttt{s\_IDLE}};
    \node[state, right=of idle] (start) {\texttt{s\_START}};
    \node[state, below=of start] (data) {\texttt{s\_DATA}};
    \node[state, left=of data] (stop) {\texttt{s\_STOP}};

    \path[->, thick]
        (idle) edge[bend left=12] node {\texttt{tx\_en}} (start)
        (idle) edge[loop above] node {\texttt{tx\_en = 0}} (idle)
        (start) edge[loop right] node {\texttt{clk\_count \textless{} CLOCKS\_PER\_BIT - 1}} (start)
        (start) edge node {\texttt{clk\_count = CLOCKS\_PER\_BIT - 1}} (data)
        (data) edge[loop below] node {\shortstack{\texttt{clk\_count \textless{} CLOCKS\_PER\_BIT - 1} \\
                 \texttt{o bit\_index \textless{} 7}}} (data)
        (data) edge node {\shortstack{\texttt{bit\_index = 7} \\
                                     \texttt{y fin de bit}}} (stop)
        (stop) edge[loop left] node {\texttt{clk\_count \textless{} CLOCKS\_PER\_BIT - 1}} (stop)
        (stop) edge[bend left=12] node {\texttt{fin de bit}} (idle);
\end{tikzpicture}
    }
\caption{Diagrama visual de la FSM del transmisor UART.}
\end{figure}

\section{Registros internos}

\begin{lstlisting}[language=Verilog]
reg [1:0]  state;
reg [15:0] clk_count;
reg [2:0]  bit_index;
reg [7:0]  tx_data_reg;
\end{lstlisting}

Cada registro tiene una funcion especifica:

\begin{itemize}
    \item \texttt{state}: guarda el estado actual de la maquina.
    \item \texttt{clk\_count}: cuenta cuantos ciclos de reloj han pasado dentro del bit actual.
    \item \texttt{bit\_index}: indica que bit del byte se esta transmitiendo, desde 0 hasta 7.
    \item \texttt{tx\_data\_reg}: almacena una copia local del byte a transmitir.
\end{itemize}

Guardar el dato en \texttt{tx\_data\_reg} es importante porque la entrada \texttt{tx\_data} podria cambiar mientras la transmision esta en curso.

\section{Bloque secuencial principal}

\begin{lstlisting}[language=Verilog]
always @(posedge clk or negedge rst_n)
\end{lstlisting}

Todo el comportamiento del transmisor ocurre en este bloque secuencial. Eso significa que:

\begin{itemize}
    \item el modulo actualiza sus registros en cada flanco positivo del reloj,
    \item y si \texttt{rst\_n} baja a \texttt{0}, se reinicia de inmediato.
\end{itemize}

\section{Comportamiento durante reset}

\begin{lstlisting}[language=Verilog]
if (!rst_n) begin
    state       <= s_IDLE;
    clk_count   <= 0;
    bit_index   <= 0;
    tx_data_reg <= 8'h00;
    tx          <= 1'b1;
    tx_busy     <= 1'b0;
end
\end{lstlisting}

Cuando ocurre reset:

\begin{itemize}
    \item el estado vuelve a \texttt{s\_IDLE},
    \item la cuenta de reloj se limpia,
    \item el indice del bit vuelve a cero,
    \item la copia local del dato se borra,
    \item la linea \texttt{tx} queda en \texttt{1},
    \item y \texttt{tx\_busy} queda en \texttt{0}.
\end{itemize}

Que \texttt{tx} quede en \texttt{1} es correcto, porque en UART la linea en reposo se mantiene alta.

\section{Estado \texttt{s\_IDLE}}

\begin{lstlisting}[language=Verilog]
s_IDLE: begin
    tx        <= 1'b1;
    clk_count <= 0;
    bit_index <= 0;
    
    if (tx_en) begin
        tx_data_reg <= tx_data;
        tx_busy     <= 1'b1;
        state       <= s_START;
    end else begin
        tx_busy     <= 1'b0;
        state       <= s_IDLE;
    end
end
\end{lstlisting}

Este es el estado de espera.

\begin{itemize}
    \item La salida \texttt{tx} se mantiene en \texttt{1}, porque no se esta transmitiendo nada.
    \item \texttt{clk\_count} y \texttt{bit\_index} se reinician para preparar una transmision nueva.
    \item Si \texttt{tx\_en} se activa, el modulo captura \texttt{tx\_data}, marca \texttt{tx\_busy = 1} y pasa a \texttt{s\_START}.
    \item Si \texttt{tx\_en} no esta activo, el modulo sigue esperando.
\end{itemize}

En otras palabras, \texttt{s\_IDLE} detecta el inicio de una nueva operacion.

\section{Estado \texttt{s\_START}}

\begin{lstlisting}[language=Verilog]
s_START: begin
    tx <= 1'b0;
    if (clk_count < CLOCKS_PER_BIT - 1) begin
        clk_count <= clk_count + 1;
        state     <= s_START;
    end else begin
        clk_count <= 0;
        state     <= s_DATA;
    end
end
\end{lstlisting}

Aqui se envia el bit de inicio.

\begin{itemize}
    \item La salida \texttt{tx} se fuerza a \texttt{0}.
    \item \texttt{clk\_count} se incrementa ciclo a ciclo.
    \item Mientras no se cumpla la duracion completa del bit, el estado no cambia.
    \item Al terminar el conteo, se reinicia \texttt{clk\_count} y se pasa a \texttt{s\_DATA}.
\end{itemize}

La idea principal es que el bit de inicio dure exactamente un tiempo de bit.

\section{Estado \texttt{s\_DATA}}

\begin{lstlisting}[language=Verilog]
s_DATA: begin
    tx <= tx_data_reg[bit_index];
    if (clk_count < CLOCKS_PER_BIT - 1) begin
        clk_count <= clk_count + 1;
        state     <= s_DATA;
    end else begin
        clk_count <= 0;
        if (bit_index < 7) begin
            bit_index <= bit_index + 1;
            state     <= s_DATA;
        end else begin
            bit_index <= 0;
            state     <= s_STOP;
        end
    end
end
\end{lstlisting}

Este es el bloque mas importante, porque aqui se transmiten los 8 bits del byte.

\subsection*{Que se envia}

La linea:

\begin{lstlisting}[language=Verilog]
tx <= tx_data_reg[bit_index];
\end{lstlisting}

hace que en la salida aparezca el bit seleccionado por \texttt{bit\_index}. Como \texttt{bit\_index} empieza en 0, el modulo envia primero el bit menos significativo (\texttt{LSB}). Ese comportamiento coincide con el protocolo UART clasico.

\subsection*{Como avanza}

\begin{itemize}
    \item Durante un tiempo de bit, \texttt{clk\_count} cuenta los ciclos del reloj.
    \item Cuando termina ese tiempo, el modulo decide si debe pasar al siguiente bit o terminar la carga util.
    \item Si \texttt{bit\_index < 7}, incrementa el indice y permanece en \texttt{s\_DATA}.
    \item Si ya se transmitio el bit 7, reinicia \texttt{bit\_index} y cambia a \texttt{s\_STOP}.
\end{itemize}

\section{Estado \texttt{s\_STOP}}

\begin{lstlisting}[language=Verilog]
s_STOP: begin
    tx <= 1'b1;
    if (clk_count < CLOCKS_PER_BIT - 1) begin
        clk_count <= clk_count + 1;
        state     <= s_STOP;
    end else begin
        clk_count <= 0;
        state     <= s_IDLE;
    end
end
\end{lstlisting}

En este estado se envia el bit de parada:

\begin{itemize}
    \item la linea \texttt{tx} vuelve a \texttt{1},
    \item se espera un tiempo de bit completo,
    \item y despues el modulo regresa a \texttt{s\_IDLE}.
\end{itemize}

Cuando vuelve a \texttt{s\_IDLE}, el transmisor ya esta listo para aceptar otro byte.

\section{Secuencia completa de transmision}

Si se quiere transmitir, por ejemplo, el byte:

\[
\texttt{tx\_data} = 8'b1010\,1100
\]

entonces el orden en la linea \texttt{tx} sera:

\begin{enumerate}
    \item reposo: \texttt{1}
    \item bit de inicio: \texttt{0}
    \item bit 0
    \item bit 1
    \item bit 2
    \item bit 3
    \item bit 4
    \item bit 5
    \item bit 6
    \item bit 7
    \item bit de parada: \texttt{1}
\end{enumerate}

Como UART transmite primero el bit menos significativo, el orden real de los bits de datos sera:

\[
0,\ 0,\ 1,\ 1,\ 0,\ 1,\ 0,\ 1
\]

\section{Papel de \texttt{tx\_busy}}

La senal \texttt{tx\_busy} se activa cuando el modulo detecta \texttt{tx\_en} en \texttt{s\_IDLE}. Su objetivo es informar que el transmisor esta ocupado y no deberia recibir una nueva peticion.

En este diseno, \texttt{tx\_busy} se limpia al volver a \texttt{s\_IDLE}. Eso hace que sea una bandera sencilla de control para el resto del sistema.

\section{Resumen funcional}

El funcionamiento general puede resumirse asi:

\begin{enumerate}
    \item el modulo espera en \texttt{s\_IDLE},
    \item cuando \texttt{tx\_en} se activa, guarda el byte,
    \item envia un bit de inicio en \texttt{s\_START},
    \item envia los 8 bits del dato en \texttt{s\_DATA},
    \item envia un bit de parada en \texttt{s\_STOP},
    \item y vuelve a reposo.
\end{enumerate}

\section{Observaciones importantes para entender el codigo}

\begin{itemize}
    \item El diseno usa una sola maquina de estados, por eso es facil seguir su flujo.
    \item La temporizacion de la UART depende directamente de \texttt{CLOCKS\_PER\_BIT}.
    \item El envio de datos es \texttt{LSB first}, como es habitual en UART.
    \item No hay bit de paridad ni multiples bits de parada; es una implementacion basica de tipo 8N1.
    \item El modulo asume que \texttt{tx\_en} se activa cuando el transmisor esta libre.
\end{itemize}

\section{Conclusion}

La implementacion de \texttt{uart\_tx.v} es un transmisor UART compacto y claro. Internamente combina:

\begin{itemize}
    \item una maquina de estados para organizar la secuencia de envio,
    \item un contador para fijar la duracion de cada bit,
    \item un indice para recorrer los 8 bits del byte,
    \item y un registro auxiliar para conservar el dato durante toda la transmision.
\end{itemize}

Desde el punto de vista de aprendizaje, este modulo es una buena base para entender comunicacion serial, temporizacion por baud rate y diseño secuencial en Verilog.

\end{document}
